% A sample LaTeX draft for the guidelines checks. A fixture, not a paper: it loads a conference
% style package so the template check has something to compare against, carries a compiled PDF
% beside it so the page count can be read without a compiler, and cites author-date so a venue
% that wants numbered citations has something to complain about.
\documentclass{article}
\usepackage{iclr2027_conference}
\usepackage{graphicx}
\usepackage{natbib}

\title{A sample LaTeX draft for the guidelines checks}

\begin{document}
\maketitle

\begin{abstract}
A short abstract, present so that the section check can find one and the word count can leave it
out of the main text.
\end{abstract}

\section{Introduction}
Venue guidelines differ, and the mechanical half of conforming to them is a page count, a section
list, a style package and a citation form \citep{sample2026}.

\section{Method}
Nothing is measured here. The draft exists so that the checks have LaTeX to read: a figure
environment with a caption, a bibliography, and a preamble that loads a conference style.

\begin{figure}
\includegraphics{figures/fig1.png}
\caption{A grey square at 300 dpi, standing in for a panel.}
\label{fig:one}
\end{figure}

\section{Results}
The draft compiles to three pages, which is what \texttt{draft.pdf} beside it says
\citep{sample2026,other2026}.

\section{Reproducibility Statement}
Everything in this fixture is deterministic.

\section{AI-use statement}
This fixture was written by a model under a ticket, which is the sort of thing this statement is
for.

\bibliographystyle{plain}
\bibliography{references}

\end{document}
